Un supermarché dispose d'un stock de tomates provenant de deux fournisseurs A et B.

II a été constaté que :

\begin{itemize}[label=$\bullet~$]
\item 91\,\% du stock de tomates est commercialisable ;
\item 60\,\% du stock de tomates provient du fournisseur A ;
\item parmi les tomates provenant du fournisseur A , la proportion de tomates commercialisables
est de 95\,\%.
\end{itemize}

On choisit au hasard une tomate dans le stock.

On désigne par :
\begin{itemize}[label=$\bullet~$]
\item $A$ l'évènement \og La tomate provient du fournisseur A \fg ;
\item $B$ l'évènement \og La tomate provient du fournisseur B \fg ;
\item $C$ l'évènement \og La tomate est commercialisable \fg.
\end{itemize}

Pour un évènement quelconque $E$, on note $P(E)$ la probabilité de $E$.

\medskip

\textbf{Partie A}

\medskip

\begin{enumerate}
\item Recopier l'arbre ci-dessous en complétant les pointillés.


\begin{center}
\pstree[treemode=R,levelsep=3.5cm, treesep=0.7cm]{\TR{}}
{\pstree{\TR{$A~$}\ncput*{\ldots}}
	{\TR{$C$} \ncput*{\ldots}
	\TR{$\overline{C}$} \ncput*{\ldots}
	}
\pstree{\TR{$B~$}\ncput*{\ldots}}
	{\TR{$C$} 
	\TR{$\overline{C}$} 
	}
}
\end{center}






















\item
	\begin{enumerate}
		\item Déterminer la probabilité que la tomate choisie soit commercialisable et provienne du fournisseur A
		\item Démontrer que $P_B(C) = 0,85$.
		\item La tomate choisie est non commercialisable. Le responsable des achats estime qu'il
y a deux fois moins de chance qu'elle provienne du fournisseur A que du fournisseur B. A-t-il raison ?
	\end{enumerate}
\end{enumerate}

\bigskip

\textbf{Partie B}

\medskip

On rappelle que 9\,\% des tomates du stock ne sont pas commercialisables.

\medskip

\begin{enumerate}
\item On prend $15$ tomates dans le stock au hasard et de manière indépendante. On considère que le stock est suffisamment important pour qu'on puisse assimiler ce prélèvement à un tirage aléatoire avec remise.

On note $X$ la variable aléatoire égale au nombre de tomates non commercialisables dans cet échantillon de $15$~tomates.
	\begin{enumerate}
		\item On admet que la variable aléatoire $X$ suit une loi binomiale. En préciser les paramètres.
		\item Déterminer la probabilité qu'exactement deux tomates soient non commercialisables.

\emph{On donnera la valeur arrondie au millième.}
		\item Déterminer la probabilité qu'au plus deux tomates soient non commercialisables.

\emph{On donnera la valeur arrondie au millième.}
	\end{enumerate}
\item On constitue désormais un échantillon de $n$ tomates, toujours dans les mêmes conditions, où $n$ désigne un entier naturel non nul.

On note $X_n$, la variable aléatoire égale au nombre de tomates non commercialisables et $F_n$ la variable aléatoire égale à la fréquence de tomates non commercialisables dans cet échantillon de $n$ tomates.

On a donc $F_n = \dfrac{X_n}{n}$.

On admet que la variable aléatoire $X_n$, suit la loi binomiale de paramètres $n$ et $0,09$.
	\begin{enumerate}
		\item Calculer l'espérance $E(F_n)$ et exprimer la variance $V(F_n)$ en fonction de $n$.
		\item Démontrer que $P(0,04 < F_n < 0,14) \geqslant 1 - \dfrac{32,76}{n}$.
		\item Le responsable des achats prélève dans le stock un échantillon de $500$ tomates. Il s'aperçoit que $55$ tomates ne sont pas commercialisables. 

Est-ce conforme à ce qu'il pouvait attendre ? Justifier la réponse.
	\end{enumerate}
\end{enumerate}

\medskip
